\documentclass[10pt]{article}
\usepackage[margin=0.9in]{geometry}
\usepackage{booktabs}
\usepackage{xcolor}
\setlength{\parskip}{4pt}
\setlength{\parindent}{0pt}
\newcommand{\todo}[1]{\textcolor{red}{[TODO: #1]}}

\title{\vspace{-1.5em}CS 3220 Lab 2 Report: Branch Predictor Experiments}
\author{\todo{name}}
\date{}

\begin{document}
\maketitle
\vspace{-2em}

\section*{Setup}

The predictor has three parts: a branch history register (BHR), a pattern
history table (PHT) of 2-bit counters, and a direct-mapped branch target
buffer (BTB). Tasks 8 and 9 swap in a 1-bit PHT and a two-way BTB. Fetch predicts taken only when the BTB hits \emph{and} the PHT
counter says taken. Otherwise it predicts PC+4. The BTB is written every time
a branch or jump resolves, and it always stores the taken target.

Accuracy is the number of correct next-PC predictions divided by the number
of resolved branches and jumps, as printed by the supplied simulator.

\section*{Task 7: Predictor sizes}

All runs use \texttt{towers.mem} (99{,}635 branches and jumps). The PHT index
is PC XOR an 8-bit BHR in every run. Only the table sizes change.

\begin{center}
\begin{tabular}{lrrrr}
\toprule
Config & BTB entries & PHT entries & Accuracy & CPI \\
\midrule
Baseline      & 16  & 256  & 33.85\% & 1.68 \\
PHT alone     & 16  & 1024 & 33.85\% & 1.68 \\
BTB alone     & 256 & 256  & 64.44\% & 1.38 \\
BTB and PHT   & 256 & 1024 & 98.40\% & 1.05 \\
\bottomrule
\end{tabular}
\end{center}

\textbf{PHT alone} does not help. Without a BTB hit, fetch has no target, so
a better direction guess cannot be used.

\textbf{BTB alone} nearly doubles the accuracy. With 16 entries the branches
in Towers keep replacing each other in the BTB. With 256 entries they stop.

\textbf{BTB and PHT} together reach 98\%. Once the BTB hits, the direction
guess is what limits accuracy. With 256 counters, different branches share
counters and undo each other's training. With 1024 counters they mostly do
not.

The submitted code uses the ``BTB and PHT'' configuration (98.40\%), which is
above the 50\% target.

\section*{Task 8: One-bit versus two-bit counters}

Both runs use \texttt{counter.mem} with 256 PC-indexed PHT entries and a
64-entry direct-mapped BTB. Only the counter width changes.

\begin{center}
\begin{tabular}{lrrr}
\toprule
Counter & Branches & Correct & Accuracy \\
\midrule
2-bit saturating & 8{,}411 & 6{,}192 & 73.62\% \\
1-bit            & 8{,}411 & 4{,}116 & 48.94\% \\
\bottomrule
\end{tabular}
\end{center}

The benchmark is a short inner loop whose branch goes taken, taken, not
taken, over and over, inside an outer loop whose branch is taken 31 times out
of 32.

\textbf{2-bit} counters get the inner branch right two times out of three.
The counter sits at ``taken''. The one not-taken outcome moves it down a
single step, which is not enough to flip the prediction, so only the loop
exit is missed.

\textbf{1-bit} counters get the inner branch right one time out of three. The
bit always holds the last outcome, so the loop exit is missed, and then the
first iteration of the next pass is missed too, because the bit now says
``not taken''. Each loop exit costs two mispredictions instead of one.

The outer branch is predicted almost perfectly in both runs. Per call that
gives about $64 + 31$ correct out of 128 for 2-bit (74\%) and $32 + 31$ out
of 128 for 1-bit (49\%), which matches the measurements.

\section*{Task 9: Direct-mapped versus two-way BTB}

Both runs use \texttt{btb.mem} with a 16-entry BTB. Only the BTB organisation
changes. The direction predictor is the same in both: 1024 two-bit counters
indexed by PC XOR an 8-bit BHR.

\begin{center}
\begin{tabular}{lrrr}
\toprule
BTB & Branches & Correct & Accuracy \\
\midrule
Direct-mapped, 16 entries   & 16{,}475 & 176 & 1.07\% \\
Two-way, 16 entries (8 sets) & 16{,}475 & 16{,}166 & 98.12\% \\
\bottomrule
\end{tabular}
\end{center}

The benchmark alternates between two taken branches that are 64 bytes apart.

\textbf{Direct-mapped.} The index is PC[5:2], so two PCs that differ by 64
bytes get the same index and only one of them fits. Each branch overwrites
the other when it resolves, so every lookup misses and fetch predicts PC+4,
which is wrong for a taken branch. Almost every prediction fails. CPI is 3.10.

\textbf{Two-way.} There are 8 sets indexed by PC[4:2], and each set holds two
tags. The two branches still land in the same set, but they sit in different
ways and stop evicting each other. Accuracy goes to 98\% and CPI drops to
1.96, with the same 16 entries of storage.

The replacement policy matters. Each set has one LRU bit that points at the
way written least recently. An earlier version always replaced way 0 when
the set was full. A branch from the start-up code stayed in way 1 for the
whole run, the two loop branches fought over way 0, and accuracy was about
1\%, no better than direct-mapped.

\end{document}
